% ch8-orders.tex
\chapter{Choosing an Evaluation Order}
\label{ch:orders}

\begin{answer}
ulm.lam's cost model counts floating-point operations (flops). With static
sizes the dynamic program of kernel selection finds one cheapest plan at
compile time. With run-time sizes it cannot, so it runs the dynamic program
for several sample sizes, keeps every plan that wins for some sample, and
emits all of them behind an \code{scf.if} that compares their flop counts,
computed from the actual sizes when the function is called. The comparison
is exact for the plans it has; whether those plans include the best one for
every size is a heuristic, not a proof.
\end{answer}

% ---------------------------------------------------------------------------
\section{The cost model}
\label{sec:orders-cost}

Costs are the flop counts in the table of \cref{sec:kernels-select}: \(2mkn\)
for a product, \(m^2n\) for a triangular solve, \(n^2 r\) for a Gram
product. They are computed inline by the dynamic program; the ops'
\code{KernelCostOpInterface} methods (\code{estimateFlops},
\code{estimateBytes}) exist but are not called
(\cref{sec:infra-analogy}). A dynamic dimension that has no sample value
counts as \(1000\) (\code{kNominal}, \srcref{lib/UlmLamPasses.cpp}{190}{200}).

Flops are a crude model: \code{gemv} moves as many bytes as it computes
flops, \code{gemm} reuses each byte many times, and \code{potrf} runs at a
lower rate than \code{gemm} for small matrices. ulm.lam contains one
measured correction: the choice between \(F P F\T\) by \code{gemm} plus
\code{gemmtr} and by Cholesky of \(P\), \code{trmm} and \code{syrk} uses a
constant \(n_0 = 2400\) fitted to OpenBLAS on an Apple M5
(\code{choleskySandwich}, \srcref{lib/UlmLamPasses.cpp}{1592}{1724}). The
README lists a measured cost model, in the spirit of ATLAS, as open work.

% ---------------------------------------------------------------------------
\section{Several plans and a run-time choice}
\label{sec:orders-runtime}

\begin{steps}
\item[Input] \(QK\T V\) with \(Q,K\in\R^{n\times d}\), \(V\in\R^{n\times e}\), all sizes known only at run time (\code{test/select/runtime_order.mlir}).
\item[Representation] one n-ary \code{ulm_lam.mul}; sizes as \code{index} parameters via \code{ulm_lam.dims}.
\item[Rule] run the dynamic program for every assignment of \(\{10, 1000\}\) to the (at most four) distinct sizes; collect the distinct optimal plans; emit their costs as \code{index} arithmetic and select with \code{scf.if}.
\item[Output] two plans, chosen at run time (below).
\item[Theory] symbolic parameters, \FG{6.1}{Symbolic compile-time parameters}; multi-versioning as a residual of partial evaluation, \FG{8.5}{Partial evaluation and the Futamura projections}; cost models, \FG{10.5}{Auto-tuning and learned cost models}.
\end{steps}

The output after selection and \code{canonicalize}
(\code{outputs/runtime_order.mlir}; \code{\%arg0} \(=n\), \code{\%arg1}
\(=d\), \code{\%arg2} \(=e\)):

\begin{lstlisting}
%0 = arith.muli %arg0, %arg1 : index      // n d
%1 = arith.muli %0, %arg0 : index         // n d n
%2 = arith.muli %1, %c2 : index           // 2 n d n        : Q K^T
%3 = arith.muli %arg0, %arg0 : index
%4 = arith.muli %3, %arg2 : index
%5 = arith.muli %4, %c2 : index           // 2 n n e        : (Q K^T) V
%6 = arith.addi %2, %5 : index            // cost of plan 0
...                                       // 2 d n e + 2 n d e: cost of plan 1
%13 = arith.addi %9, %12 : index
%14 = tensor.empty(%arg0, %arg2) : tensor<?x?xf64>
%15 = arith.cmpi ule, %6, %13 : index
%16 = scf.if %15 -> (tensor<?x?xf64>) {
  %17 = tensor.empty(%arg0, %arg0) : tensor<?x?xf64>
  %18 = ulm_lamk.gemm trans_a(none) trans_b(trans) ... ins(%arg3, %arg4 ...) outs(%17 ...)   // Q K^T
  %19 = ulm_lamk.gemm trans_a(none) trans_b(none) ... ins(%18, %arg5 ...) outs(%14 ...)    // (Q K^T) V
  scf.yield %19 : tensor<?x?xf64>
} else {
  %17 = tensor.empty(%arg1, %arg2) : tensor<?x?xf64>
  %18 = ulm_lamk.gemm trans_a(trans) trans_b(none) ... ins(%arg4, %arg5 ...) outs(%17 ...)  // K^T V
  %19 = ulm_lamk.gemm trans_a(none) trans_b(none) ... ins(%arg3, %18 ...) outs(%14 ...)    // Q (K^T V)
  scf.yield %19 : tensor<?x?xf64>
}
\end{lstlisting}

Both branches write into the same destination \code{\%14}, so that
bufferization gives the \code{scf.if} result one buffer instead of one per
branch (\srcref{lib/UlmLamPasses.cpp}{1847}{1907}). The README measures the
effect: for \(n=4096\), \(d=64\), the formula as written in Julia,
\((QK\T)V\), takes 43 times as long as ulm.lam, which picks \(Q(K\T V)\).

\paragraph{How the plans are found} \code{plansOf}
(\srcref{lib/UlmLamPasses.cpp}{1495}{1528}) samples the corners of the size
space, \(\{10, 1000\}^s\) for \(s \le 4\) distinct sizes, and keeps each
distinct optimal plan. For the chain above two plans suffice for all sizes,
but in general a plan that is optimal only in the interior of the size
space would be missed, and with more than four distinct sizes only the
nominal plan is generated.

\begin{theory}[ (\FGs{6.1}, \FGchs{6})]
The costs are polynomials in symbolic sizes. Deciding at compile time which
plan is cheapest for all admissible sizes is a question about the sign of a
polynomial difference over a region: the kind of symbolic reasoning of
\FGch{6}{Constraints and Symbolic Reasoning}, with nonlinear (polynomial)
constraints that Presburger arithmetic does not cover. ulm.lam sidesteps the
decision: it samples to find candidate plans and defers the comparison to
run time, when the sizes are numbers. The values are unknown at compile
time but become known before the computation starts, which is the binding
time situation of \FG{6.1}{Symbolic compile-time parameters}.
\end{theory}

\begin{theory}[ (\FGs{8.5}, \FGs{9.2}, \FGs{10.5})]
A function with several specialized versions and a dispatch on run-time
values is a residual program in the sense of
\FG{8.5}{Partial evaluation and the Futamura projections}: everything that
does not depend on the sizes was decided at compile time. Unlike a JIT
(just-in-time) compiler with speculation, \FG{9.2}{JIT compilation:
tiering, OSR, deoptimization}, nothing is guessed and nothing needs to be
undone; every version is correct for every size. And the cost model is an
analytical one, not learned or measured, \FG{10.5}{Auto-tuning and learned
cost models}.
\end{theory}

% ---------------------------------------------------------------------------
\section{Decisions that depend on a trip count}
\label{sec:orders-hoist}

In a loop, \(b_i = S\inv H\T R\inv y_i\) can be evaluated as written in
every step, or the operator \(P = S\inv H\T R\inv\) can be formed once before
the loop and applied with one \code{gemv} per step. The second pays off if
the loop runs often enough:
\[
\mathrm{cost}(P) + \mathrm{trips}\cdot 2\,\mathrm{rows}(P)\,\mathrm{cols}(P)
\;<\; \mathrm{trips}\cdot\mathrm{cost}(\text{step}).
\]
The trip count is usually known only at run time, so \code{hoistOperator}
(\srcref{lib/UlmLamPasses.cpp}{1909}{2064}) emits the condition and, if it
holds, \(P\) before the loop, and both variants in the loop under the same
condition. Matrix-vector flops are weighted ten times matrix-matrix flops
(\code{kLevel2 = 10}) in this decision, a second crude correction of the flop model for
memory-bound kernels. The Kalman-filter examples
(\code{examples/kalman/}) exercise it.

% ---------------------------------------------------------------------------
\section{Specialization by element type, not by size}
\label{sec:orders-spec}

One compiled library serves all sizes. What the Julia frontend does
specialize is the element type: a call with \code{Float32} arrays compiles
a second library from the same traced text with \code{f64} replaced by
\code{f32} (\srcref{julia/UlmLam/src/runtime.jl}{262}{264}). Libraries are
cached under a hash of the program text, the compiler's modification time,
the backend and the flags.

\begin{theory}[ (\FGs{8.7})]
Compiling one copy per element type and choosing it at the call by the
argument types is monomorphization, \FG{8.7}{Monomorphization vs dictionary
passing}. Sizes are handled the other way: one generic code with run-time
size parameters, plus the run-time plan selection of
\cref{sec:orders-runtime} where the order depends on them.
\end{theory}
